\documentclass[a4paper,10pt]{article}

\usepackage[T1]{fontenc}
\usepackage[utf8]{inputenc}
\usepackage{lmodern}
\usepackage[margin=1in]{geometry}
\usepackage{amsmath,amssymb,amsthm}
\usepackage{booktabs}
\usepackage{graphicx}
\usepackage{microtype}
\usepackage{pdflscape}
\usepackage[hidelinks]{hyperref}
\hypersetup{pdftitle={Supplementary Information for The Architecture of Insight: A Reduced-Model Hypothesis for Representational Flexibility and Insight Capacity}}

\pagestyle{plain}
\setlength{\parindent}{0pt}

\begin{document}

\begin{center}
{\Large\bfseries Supplementary Information}
\end{center}

\section*{Supplementary Note S1. Transverse minimization and slaving}

Let $v^*$ be the normalized contrastive eigenvector, let
$V_\perp=\{w:w^Tv^*=0\}$, and write $S=xv^*+y$ with $y\in V_\perp$. For fixed
$\varepsilon$, define
\[
A:=\bigl(I-W(\varepsilon)\bigr)\big|_{V_\perp},\qquad
q:=P_\perp (v^*)^{\circ3},
\]
where $P_\perp$ is the orthogonal projector onto $V_\perp$ and the superscript
$\circ3$ denotes componentwise cubing. Hypothesis (A2) gives $A\succeq\gamma I$.

For the continuous Hopfield energy
\[
E(S)=-\frac12S^TW(\varepsilon)S+\sum_i F(S_i),\qquad
F'(z)=\tanh^{-1}(z),
\]
the Hessian is $H(S)=-W(\varepsilon)+D(S)$, with
$D(S)=\operatorname{diag}((1-S_i^2)^{-1})\succeq I$. Therefore, for every
$w\in V_\perp$ inside the activation cube,
\[
w^TH(S)w\geq\gamma\lVert w\rVert^2.
\]
The energy is thus uniformly strongly convex on each transverse slice, so that
the slice has at most one minimizer. At $(x,y)=(0,0)$ the transverse stationarity
equation has derivative $A$, and the implicit function theorem supplies a unique
smooth branch $y^*(x)$ locally. Because $E(-S)=E(S)$, uniqueness implies
$y^*(-x)=-y^*(x)$.

Projecting the stationarity equation onto $V_\perp$ and using
$W(\varepsilon)v^*=\lambda_W(\varepsilon)v^*$ gives
\[
0=P_\perp\nabla E(xv^*+y)
=Ay+\frac{x^3}{3}q+O(x^2\lVert y\rVert+\lVert y\rVert^3+x^5).
\]
Consequently,
\[
\boxed{y^*(x)=-\frac{x^3}{3}A^{-1}q+O(x^5)},
\qquad
\lVert y^*(x)\rVert
\leq\frac{|x|^3\sqrt{h-g^2}}{3\gamma}+O(x^5),
\]
where $g=\sum_i(v_i^*)^4$, $h=\sum_i(v_i^*)^6$, and
$\lVert q\rVert^2=h-g^2$. This proves the slaving statement used in the main
text. It also shows explicitly why transverse minimization first affects the
reduced energy at order $x^6$.

\section*{Supplementary Note S2. Sextic coefficient and normal-form barrier}

On the line $S=xv^*$, the saturation expansion is
\[
\sum_i F(xv_i^*)
=\frac{x^2}{2}+\frac{g}{12}x^4+\frac{h}{30}x^6+O(x^8).
\]
Combining this with the bilinear term and
$\mu(\varepsilon)=\lambda_W(\varepsilon)-1$ gives
\[
E(xv^*)=-\frac{\mu}{2}x^2+\frac{g}{12}x^4+\frac{h}{30}x^6+O(x^8).
\]
For $y=O(x^3)$, the terms needed through order $x^6$ are
\[
E(xv^*+y)=E(xv^*)+\frac{x^3}{3}y^Tq+\frac12y^TAy+O(x^8).
\]
Minimization in $y$ inserts $y=-(x^3/3)A^{-1}q+O(x^5)$ and yields
\[
U(x;\varepsilon)
=-\frac{\mu}{2}x^2+\frac{g}{12}x^4+c_6x^6+O(x^8),
\qquad
\boxed{c_6=\frac{h}{30}-\frac1{18}q^TA^{-1}q}.
\]
Since $0\preceq A^{-1}\preceq\gamma^{-1}I$ on $V_\perp$,
\[
\frac{h}{30}-\frac{h-g^2}{18\gamma}
\leq c_6\leq\frac{h}{30}.
\]
The Level-B model used for the dynamical analysis discards the sextic and higher
terms:
\[
U_{\mathrm{nf}}(x;\varepsilon)=\frac{g}{12}x^4-\frac{\mu}{2}x^2.
\]
Its wells and barrier are
\[
x_\pm=\pm\sqrt{\frac{3\mu}{g}},\qquad
\Delta E_{\mathrm{nf}}=\frac{3\mu^2}{4g}.
\]
Treating $c_6x^6$ perturbatively gives the leading relative correction
\[
\frac{\Delta E-\Delta E_{\mathrm{nf}}}{\Delta E_{\mathrm{nf}}}
=-\frac{36c_6}{g^2}\mu+O(\mu^2).
\]
These statements concern the near-bifurcation asymptotic energy reduction. The
one-dimensional stochastic dynamics used in the production sweep additionally
assume adiabatic transverse slaving and do not quantify noise-induced
renormalization from transverse fluctuations.

\section*{Supplementary verification note}

The supplied deterministic symbolic and numerical audit reproduced all eight
reported sanity checks without differences from its archived output. The checks
covered the saturation series, normal-form stationary points and barrier,
canonical rescaling, the actions of the Hebbian perturbation on $v^*$ and $u$,
the $x^6$ slaving coefficient, transverse convexity, and a counterexample in
which (A2) fails. These are consistency checks, not formal verification; no proof
assistant was used.

\clearpage
\begin{landscape}

\begin{center}
{\large\bfseries Supplementary Table S1. Complete production sweep}
\end{center}

\vspace{0.5em}

\small
\setlength{\tabcolsep}{4pt}
\renewcommand{\arraystretch}{1.12}
\begin{center}
\resizebox{\textwidth}{!}{%
\begin{tabular}{@{}rrrrrrrrrrrrrrrr@{}}
\toprule
$r$ & Seed & $n$ & Events &
$\kappa_{\rm sim}$ & CI low & CI high &
$\kappa_{\rm exact}$ & $\kappa_K$ &
Err. exact & Err. $K$ &
$\log(S/E)$ & $\log(S/K)$ & $\log(E/K)$ &
Est. diff. (\%) & $t^*$\\
\midrule
0.1500 & 77000 & 6000 & 6000 & 0.51535 & 0.44072 & 0.60261 & 0.44985 & 0.19373 & 0.1456 & 1.6602 & 0.1359 & 0.9784 & 0.8425 & 22.856 & 8.00\\
0.1938 & 87000 & 6000 & 6000 & 0.39050 & 0.38043 & 0.40084 & 0.35971 & 0.18542 & 0.0856 & 1.1060 & 0.0821 & 0.7448 & 0.6626 & 12.990 & 0.50\\
0.2504 & 97000 & 6000 & 6000 & 0.31389 & 0.30541 & 0.32261 & 0.28629 & 0.17522 & 0.0964 & 0.7914 & 0.0920 & 0.5830 & 0.4910 & 15.080 & 1.00\\
0.3235 & 107000 & 6000 & 6000 & 0.23487 & 0.22894 & 0.24094 & 0.22653 & 0.16287 & 0.0368 & 0.4421 & 0.0361 & 0.3661 & 0.3299 & 9.847 & 0.50\\
0.4180 & 117000 & 6000 & 6000 & 0.18619 & 0.18153 & 0.19098 & 0.17788 & 0.14819 & 0.0468 & 0.2565 & 0.0457 & 0.2283 & 0.1826 & 8.615 & 0.50\\
0.5400 & 127000 & 6000 & 6000 & 0.14068 & 0.13716 & 0.14429 & 0.13823 & 0.13116 & 0.0177 & 0.0726 & 0.0176 & 0.0701 & 0.0525 & 6.832 & 0.50\\
0.6977 & 137000 & 6000 & 6000 & 0.11418 & 0.11130 & 0.11713 & 0.10586 & 0.11202 & 0.0786 & 0.0192 & 0.0756 & 0.0191 & $-0.0566$ & 9.850 & 1.00\\
0.9015 & 147000 & 6000 & 6000 & 0.08287 & 0.08079 & 0.08500 & 0.07936 & 0.09137 & 0.0442 & 0.0931 & 0.0433 & $-0.0977$ & $-0.1410$ & 7.696 & 1.00\\
1.1647 & 157000 & 6000 & 5999 & 0.06265 & 0.06106 & 0.06427 & 0.05764 & 0.07023 & 0.0869 & 0.1079 & 0.0833 & $-0.1142$ & $-0.1975$ & 9.838 & 2.00\\
1.5049 & 167000 & 6000 & 5992 & 0.04246 & 0.04139 & 0.04356 & 0.03991 & 0.04998 & 0.0640 & 0.1504 & 0.0620 & $-0.1630$ & $-0.2250$ & 7.234 & 2.00\\
1.9443 & 177000 & 6000 & 5903 & 0.02664 & 0.02597 & 0.02733 & 0.02573 & 0.03221 & 0.0357 & 0.1727 & 0.0350 & $-0.1896$ & $-0.2246$ & 2.691 & 1.00\\
2.5121 & 187000 & 6000 & 5400 & 0.01506 & 0.01466 & 0.01547 & 0.01492 & 0.01825 & 0.0094 & 0.1751 & 0.0094 & $-0.1925$ & $-0.2018$ & 1.673 & 1.00\\
3.2456 & 197000 & 6000 & 4044 & 0.00739 & 0.00717 & 0.00762 & 0.00743 & 0.00877 & 0.0051 & 0.1568 & $-0.0052$ & $-0.1706$ & $-0.1654$ & 0.000 & 0.00\\
4.1934 & 207000 & 6000 & 2144 & 0.00293 & 0.00281 & 0.00306 & 0.00300 & 0.00340 & 0.0222 & 0.1378 & $-0.0225$ & $-0.1483$ & $-0.1258$ & 0.000 & 0.00\\
5.4179 & 217000 & 6000 & 788 & 0.00094 & 0.00087 & 0.00100 & 0.00091 & 0.00100 & 0.0263 & 0.0633 & 0.0259 & $-0.0654$ & $-0.0914$ & 0.000 & 0.00\\
7.0000 & 227000 & 6000 & 171 & 0.000193 & 0.000166 & 0.000225 & 0.000192 & 0.000205 & 0.0061 & 0.0578 & 0.0061 & $-0.0596$ & $-0.0656$ & 0.339 & 0.50\\
\bottomrule
\end{tabular}
}%
\end{center}

\vspace{0.5em}
\footnotesize
\textit{Notes.} $S$, $E$, and $K$ denote stabilized simulation, numerically evaluated exact
one-dimensional MFPT rate, and Kramers asymptotic rate. The relative differences in the
\textup{Err.} columns are absolute. The estimator difference is $100\,|S-S_0|/S_0$, where
$S_0$ is the simulation hazard estimate from $t=0$ without truncation. Confidence intervals
use the post-truncation event count selected by the same hazard-stabilization procedure as the
central estimate. Values are rounded for display; machine-precision values are available in
\texttt{data/final\_sweep\_data\_CORRECTED\_CI.csv}.

\end{landscape}

\end{document}
